% ============================================================
% sec/04metodologia.tex
%
% Formalización con formulación matemática, pipeline previsto,
% protocolo experimental con particiones, validación y métricas
% justificadas, plan de baselines, y riesgos con mitigación.
% ============================================================
\section{Metodología Propuesta}
\label{sec:metodologia}

Esta sección describe lo que se propone hacer. Ninguna cifra de desempeño
aparece aquí: no se ha entrenado ningún modelo, y las únicas magnitudes
medidas que se citan son propiedades del corpus, establecidas por la
auditoría que documenta la Sección~\ref{sec:analisistecnico}.

\subsection{Formalización del problema}
\label{subsec:formalizacion}

Sea $\mathcal{E}$ el conjunto de partidas disponibles. Una partida
$e \in \mathcal{E}$ produce una secuencia ordenada de estados; de ella se
retiene la subsecuencia correspondiente al jugador que actúa en cada
momento. Cada estado retenido genera una observación
\[
  (x_i,\; y_i,\; g_i,\; \tau_i,\; d_i),
\]
donde $x_i \in \mathcal{X} \subseteq \mathbb{R}^p$ es el vector de
características construido con la información disponible para ese jugador en
ese instante, $y_i \in \{0,1\}$ indica si venció, $g_i$ identifica la partida
de la que procede la observación, $\tau_i$ es el número de turno y $d_i$ la
fecha. Los tres últimos se conservan para agrupar, estratificar y reportar,
pero nunca forman parte de $x_i$.

La tarea es aprender una función $f : \mathcal{X} \rightarrow [0,1]$ tal que
\begin{equation}
\label{eq:objetivo}
  f(x) \approx \Pr\!\left(Y = 1 \mid X = x\right),
\end{equation}
donde la probabilidad se toma sobre la aleatoriedad del juego posterior al
instante observado: el orden restante del mazo, la información oculta del
oponente y las decisiones futuras de ambos jugadores. Esa aleatoriedad es lo
que distingue este problema de uno determinista y lo que obliga a estimar
una probabilidad en lugar de deducir una etiqueta.

\subsubsection{La función objetivo, y qué significa cada término}

Puesto que la salida es una probabilidad, el criterio de ajuste no puede ser
una pérdida de clasificación. Se emplea la log-verosimilitud negativa con
regularización:

\begin{equation}
\label{eq:ajuste}
\begin{split}
f^{*} = \arg\min_{f \in \mathcal{F}} \Big[
  &-\frac{1}{N}\sum_{i=1}^{N}
   \Big( y_i \log f(x_i) \\[-2pt]
  &\qquad + (1-y_i)\log\big(1-f(x_i)\big) \Big)
  \;+\; \lambda\,\Omega(f) \Big].
\end{split}
\end{equation}

En lenguaje llano, la expresión busca la función que produce las
probabilidades más coherentes con lo que efectivamente ocurrió, penalizando
a la vez que esa función sea innecesariamente complicada. Los elementos son
los siguientes.

\begin{itemize}
  \item $x_i$ es el vector de características observables del estado $i$:
        premios restantes, puntos de vida, energías, cartas en mano y demás
        campos del diccionario descrito en la
        Sección~\ref{subsec:diccionario}.
  \item $y_i$ es el desenlace de la partida a la que pertenece ese estado,
        con $1$ para victoria y $0$ para derrota, desde la perspectiva del
        jugador que actúa.
  \item $f(x_i)$ es la probabilidad de victoria que el modelo asigna a ese
        estado, un número entre $0$ y $1$.
  \item $N$ es el número de observaciones empleadas en el ajuste.
  \item El primer término es la \textbf{pérdida logarítmica promedio}. Para
        cada observación toma el logaritmo de la probabilidad asignada al
        desenlace correcto. Si el modelo asigna una probabilidad alta al
        desenlace que efectivamente ocurrió, el término contribuye poco; si
        asigna una probabilidad baja, contribuye mucho. La penalización
        crece sin límite conforme la probabilidad asignada al desenlace
        correcto se acerca a cero, de modo que una predicción segura y
        equivocada resulta especialmente costosa.
  \item $\Omega(f)$ es el \textbf{término de regularización}, que penaliza la
        complejidad del modelo. Su forma concreta depende de la familia:
        para la regresión logística es la norma $\ell_2$ de los coeficientes,
        de modo que penaliza coeficientes grandes; para los modelos basados
        en árboles corresponde a los límites de profundidad, número mínimo de
        observaciones por hoja y número de hojas.
  \item $\lambda$ regula el equilibrio entre ajustarse a los datos de
        entrenamiento y mantener el modelo simple. Un valor alto produce un
        modelo más rígido; uno bajo, uno que puede memorizar. Se selecciona
        sobre el subconjunto de validación, nunca sobre el de prueba.
  \item $\mathcal{F}$ es la familia de modelos candidatos, distinta para cada
        nivel de la escalera descrita en la
        Sección~\ref{subsec:planbaselines}.
  \item $\arg\min$ indica que se elige la función $f$ dentro de
        $\mathcal{F}$ que minimiza el valor total de la expresión.
\end{itemize}

Dos precisiones sobre su papel en este trabajo. La primera es que el objetivo
no consiste en clasificar victoria o derrota, sino en producir una
probabilidad utilizable: minimizar esta expresión es lo que hace que el
óptimo teórico sea la probabilidad condicional verdadera, y por eso resulta
compatible con evaluar calibración. La segunda es que tanto $f$ como los
parámetros de preprocesamiento se ajustan exclusivamente con los datos de
entrenamiento de cada pliegue; emplear datos de validación o prueba en ese
ajuste constituiría fuga de información.

\subsubsection{Unidad de observación y unidad independiente}

Esta distinción gobierna todo el diseño experimental.

\begin{itemize}
  \item \textbf{Unidad de observación:} un estado del jugador que actúa,
        proyectado a una perspectiva común entre propio y oponente. Es la
        fila de la tabla y la entrada del modelo.
  \item \textbf{Unidad independiente:} la partida. El desenlace se sortea una
        sola vez por partida, y todas sus observaciones lo comparten.
\end{itemize}

Sobre la muestra auditada, el número de observaciones excede en
aproximadamente dos órdenes de magnitud el de partidas. Toda mención del
tamaño del conjunto en este documento expresa ambas cifras, porque reportar
solo el conteo de filas transmitiría una noción inexacta de cuánta evidencia
independiente hay disponible \cite{brill_winprob_2026}.

De ahí se sigue la restricción que rige toda partición. Para cualesquiera
dos subconjuntos $A$ y $B$ de una división,
\begin{equation}
\label{eq:agrupacion}
  \{\, g_i : i \in A \,\} \;\cap\; \{\, g_i : i \in B \,\} \;=\; \emptyset ,
\end{equation}
es decir, ninguna partida aporta observaciones a más de un subconjunto. Esta
es la formulación del argumento de tamaño muestral efectivo de Karbalaie
\textit{et al.}: la unidad independiente es la partida y no la fila
\cite{karbalaie_participantaware_2026}.

Se conservan todas las observaciones de cada partida, sin submuestreo. La
alternativa de retener una sola por partida resultaría intuitivamente
atractiva, pues produciría filas independientes, pero Brill \textit{et al.}
reportan que es peor: el conjunto completo aporta información sobre la
estructura del espacio de características que se pierde al descartar
observaciones \cite{brill_winprob_2026}.

\subsubsection{Definición del desenlace y exclusiones}

La etiqueta se deriva del registro de recompensas de cada partida, que
resuelve el desenlace en prácticamente la totalidad de los casos auditados.
Dos grupos se excluyen, con la regla declarada antes de examinar los datos
de prueba.

Los empates se excluyen del objetivo binario; su volumen es demasiado
reducido para sostener una tercera clase. Las partidas terminadas de forma
anómala, por agotamiento de tiempo, error o estado inválido, tienen
ganador identificable, pero su final no procede de la dinámica del juego, de
modo que se apartan del entrenamiento principal y se reservan para un
análisis de sensibilidad.

\subsection{Tubería propuesta}
\label{subsec:pipeline}

El procedimiento va de la repetición cruda a la tabla evaluable en siete
etapas, representadas en la Figura~\ref{fig:pipeline}.

\begin{figure}[!t]

\centering

\begin{tikzpicture}[
  font=\scriptsize,
  etapa/.style={
    draw=tecrule,
    fill=boxbg,
    thick,
    rounded corners=2pt,
    text width=5.2cm,
    align=center,
    inner sep=4pt,
    minimum height=7mm
  },
  hito/.style={
    draw=tecnavy,
    fill=tecnavy!8,
    thick,
    rounded corners=2pt,
    text width=5.2cm,
    align=center,
    inner sep=4pt,
    minimum height=7mm
  },
  flecha/.style={
    -{Stealth[length=4pt]},
    draw=tecsteel,
    thick
  },
  nota/.style={
    font=\tiny\itshape,
    text=labelgray,
    align=left,
    text width=1.9cm
  },
  fase/.style={
    font=\tiny\bfseries,
    text=tecnavy,
    align=center,
    text width=1.3cm
  }
]

\node[etapa] (e1)
  {1. Lectura y validación de la repetición};

\node[etapa, below=3.2mm of e1] (e2)
  {2. Estados del jugador activo, perspectiva normalizada};

\node[hito, below=3.2mm of e2] (e3)
  {3. Filtro de campos excluidos};

\node[etapa, below=3.2mm of e3] (e4)
  {4. Construcción de características};

\node[hito, below=3.2mm of e4] (e5)
  {5. Materialización de la tabla};

\node[hito, below=3.2mm of e5] (e6)
  {6. Partición agrupada por partida};

\node[etapa, below=3.2mm of e6] (e7)
  {7. Ajuste, selección y evaluación};

\draw[flecha] (e1) -- (e2);
\draw[flecha] (e2) -- (e3);
\draw[flecha] (e3) -- (e4);
\draw[flecha] (e4) -- (e5);
\draw[flecha] (e5) -- (e6);
\draw[flecha] (e6) -- (e7);

\node[nota, right=2.5mm of e2]
  {propio frente a oponente};

\node[nota, right=2.5mm of e3]
  {falla si entra un campo prohibido};

\node[nota, right=2.5mm of e5]
  {una sola vez, formato columnar};

\node[nota, right=2.5mm of e6]
  {precede a toda transformación};

\node[nota, right=2.5mm of e7]
  {escenarios A, B y C};

\begin{scope}[on background layer]

  \node[
    draw=tecrule,
    dashed,
    rounded corners=3pt,
    inner sep=3pt,
    fit=(e1)(e5)
  ] (faseA) {};

  \node[
    draw=tecrule,
    dashed,
    rounded corners=3pt,
    inner sep=3pt,
    fit=(e6)(e7)
  ] (faseB) {};

\end{scope}

\node[
  fase,
  left=4mm of faseA.west,
  anchor=east
] {Construcción\\del conjunto};

\node[
  fase,
  left=4mm of faseB.west,
  anchor=east
] {Experimentación};

\end{tikzpicture}

\caption{Tubería propuesta. Las etapas 1 a 5 se ejecutan una sola vez y
producen la tabla; las etapas 6 y 7 operan sobre ella. El filtro de campos
excluidos y la partición agrupada se imponen mediante pruebas automáticas y
no por convención.}

\label{fig:pipeline}

\end{figure}

Tres decisiones de esa tubería requieren justificación.

\noindent\textbf{Normalización de la perspectiva.} Cada fila se reorganiza
distinguiendo entre la información propia y la del oponente, en lugar de
identificarlos por su posición en la mesa. Así el modelo aprende a comparar
magnitudes y no a asociar el desenlace con una posición concreta. La
auditoría hace esta transformación necesaria: en los días tardíos de la
muestra, la posición en la mesa y el orden de turno coinciden casi siempre,
de modo que sin normalizar, la posición se convertiría en un atajo
disponible.

\noindent\textbf{Ajuste del preprocesamiento dentro del pliegue.} Los
parámetros de escalado, codificación de categóricas, imputación y
ponderación de clases se calculan únicamente con los datos de entrenamiento
y después se aplican al resto. Calcularlos sobre el conjunto completo antes
de particionar constituye fuga de información aunque el modelo se ajuste
después de la partición \cite{sasse_leakage_2025}.

\noindent\textbf{Materialización única.} La tabla se escribe una sola vez en
formato columnar. El costo computacional dominante es el análisis de las
repeticiones, no el entrenamiento, de modo que repetir esa lectura en cada
experimento sería el cuello de botella del proyecto.

\subsection{Protocolo experimental}
\label{subsec:protocolo}

\subsubsection{Particiones}

El protocolo primario emplea una partición $60/15/10/15$ en entrenamiento,
validación, calibración y prueba, agrupada por partida conforme a
\eqref{eq:agrupacion}. Los cuatro subconjuntos responden a funciones
distintas: el primero ajusta el modelo, el segundo selecciona
hiperparámetros, el tercero ajusta la corrección de probabilidades descrita
en la Sección~\ref{subsec:recalibracion}, y el cuarto se toca una sola vez,
al final.

El conjunto de calibración se reserva de forma explícita porque la
corrección de probabilidades debe ajustarse sobre datos que el modelo no vio
durante el entrenamiento \cite{silvafilho_calibration_2023}. Niculescu-Mizil
y Caruana detallan la consecuencia de incumplirlo: si el modelo separa
perfectamente el conjunto sobre el que se calibra, la transformación
aprendida degenera \cite{niculescumizil_probabilities_2005}.

La alternativa de particionar por fila se rechaza sobre bases medidas y no
de principio. Sobre la muestra auditada, una partición aleatoria por
observación colocaría prácticamente todas las partidas a ambos lados de la
división, de modo que el conjunto de prueba contendría situaciones casi
idénticas a las de entrenamiento. Karbalaie \textit{et al.} cuantifican el
efecto de ese error sobre datos con estructura comparable: un mismo modelo
pasa de 0.91 de acierto bajo validación estratificada a 0.66 bajo validación
que respeta la agrupación \cite{karbalaie_participantaware_2026}.

Una ventaja de este corpus respecto de los escenarios que estudian Figueiredo
y Mendes es que la agrupación no requiere inferirse: el identificador de
partida la determina de forma exacta, mientras que en su caso la
contribución central consiste precisamente en un procedimiento para deducir
agrupaciones que no vienen dadas \cite{figueiredo_leakage_2024}.

\subsubsection{Cobertura de cartas en el conjunto de entrenamiento}

Las tres particiones anteriores separan partidas, y el escenario C separa
además mazos. Ninguna de las dos separaciones debe producir cartas ausentes
del entrenamiento, y esa condición se declara aquí de forma explícita.

\begin{quote}
Toda carta presente en el corpus debe aparecer al menos una vez en el
conjunto de entrenamiento. Los conjuntos de validación, calibración y prueba
pueden contener partidas y mazos no vistos, pero sus cartas deben haber
aparecido previamente en entrenamiento.
\end{quote}

La razón es que evaluar sobre mazos no vistos y evaluar sobre cartas no
vistas son preguntas distintas. La primera examina si el modelo aprendió a
leer el estado o memorizó qué combinaciones de cartas vencen; la segunda
examina si su representación generaliza a contenido nuevo, que es lo que
Bertram \textit{et al.} estudian reservando un conjunto completo de cartas
\cite{bertram_cardrepresentations_2024}. Mezclarlas impediría interpretar
una caída de desempeño, porque no se sabría a cuál de las dos causas
atribuirla.

El procedimiento de auditoría de cobertura se ejecuta después de particionar
y antes de entrenar, y consta de tres pasos verificables:

\begin{enumerate}
  \item Construir el inventario de identificadores de carta presentes en el
        corpus completo, y el inventario correspondiente al conjunto de
        entrenamiento resultante de la partición.
  \item Calcular la diferencia entre ambos. Si no es vacía, la partición se
        rechaza y se regenera con una semilla distinta, o se reasigna al
        entrenamiento la partida de menor duración que contenga cada carta
        faltante.
  \item Registrar el inventario final, el número de cartas que aparecen una
        sola vez y cualquier carta excluida por filtros previos, de modo que
        el reporte declare la cobertura efectiva y no la supuesta.
\end{enumerate}

La evaluación sobre cartas no vistas, si llega a realizarse, se declara como
experimento independiente y con su propia partición, separado del protocolo
principal y reportado por separado.

\subsubsection{Tres escenarios de generalización}

Los tres responden preguntas distintas y se reportan por separado, nunca
promediados. La Tabla~\ref{tab:protocolos} los resume y los liga a las
dimensiones de la pregunta de investigación.

\begin{table}[!t]
\centering
\caption{Escenarios de evaluación y dimensión de la pregunta general que
responde cada uno.}
\label{tab:protocolos}
\footnotesize
\begin{tabularx}{\columnwidth}{@{}P{1.1cm} P{2.3cm} Y@{}}
\toprule
\textbf{Escenario} & \textbf{Criterio de partición} &
\textbf{Pregunta que responde}\\
\midrule
A & Partidas disjuntas, asignación aleatoria &
¿Generaliza a partidas que no vio?\\
B & Partidas disjuntas, corte cronológico &
¿Se mantiene el desempeño sobre periodos posteriores?\\
C & Partidas disjuntas, mazos disjuntos &
¿Aprendió a leer el estado o memorizó qué mazos vencen?\\
\bottomrule
\end{tabularx}
\end{table}

\noindent\textbf{Escenario B.} La partición cronológica asigna los días más
tempranos de la muestra al entrenamiento, los intermedios a validación y
calibración, y los más tardíos a prueba. El primer día se mantiene fuera por
corresponder a un arranque parcial de la competencia, y un día de
concentración anómala de mazos y agentes se reserva y se decide tras el
análisis exploratorio.

Este escenario implementa la separación entre aprender y predecir que Kaufman
\textit{et al.} describen \cite{kaufman_leakage_2011}, y responde a la
evidencia de que la población de estrategias evoluciona de forma continua
\cite{saravanan_metadiscovery_2024}. Hodge \textit{et al.} adoptan una
partición equivalente para no emplear datos posteriores en la predicción de
eventos anteriores \cite{hodge_winprediction_2021}.

\noindent\textbf{Escenario C.} Reserva mazos completos, o pares de mazos
enfrentados, para el conjunto de prueba, respetando la condición de cobertura
de cartas del apartado anterior. Es el diseño que separa la memorización de
contenido conocido de la interpretación del estado
\cite{bertram_cardrepresentations_2024,angliss_vgcbench_2026}. El criterio
exacto de separación depende del análisis de concentración de mazos que
presenta la Sección~\ref{subsec:representatividad}.

\subsubsection{Estrategia de validación}

Los hiperparámetros se seleccionan mediante validación cruzada agrupada por
partida sobre el subconjunto de entrenamiento. Emplear la variante por
defecto, que asume independencia entre observaciones, anularía el control
establecido en la partición.

Karbalaie \textit{et al.} reportan que la validación cruzada agrupada por
grupos coincide con el diseño de dejar un grupo fuera dentro de 0.02 de
acierto y a un costo un orden de magnitud menor
\cite{karbalaie_participantaware_2026}, lo que resulta decisivo cuando la
unidad de agrupación se cuenta en decenas de miles. Se prefiere un diseño
anidado allá donde el presupuesto computacional lo permita, puesto que tanto
ese trabajo como Sasse \textit{et al.} identifican el sesgo de selección como
una fuga residual cuando los mismos pliegues sirven para ajustar y para
evaluar \cite{sasse_leakage_2025}.

\subsubsection{Métricas}

Cada métrica se incluye porque responde una pregunta que las demás no pueden
responder. La Figura~\ref{fig:descomposicion} ilustra por qué hacen falta al
menos dos familias.

\begin{figure}[!t]
\centering
\begin{tikzpicture}[
  font=\scriptsize,
  caja/.style={draw=tecsteel, fill=boxbg, rounded corners=2pt,
               text width=2.6cm, align=center, inner sep=4pt,
               minimum height=8mm},
  raiz/.style={draw=tecnavy, fill=tecnavy!10, rounded corners=2pt,
               text width=3.4cm, align=center, inner sep=4pt,
               minimum height=8mm},
  met/.style={font=\tiny\itshape, text=labelgray, align=center,
              text width=2.6cm},
  flecha/.style={-{Stealth[length=4pt]}, draw=tecsteel, thick}
]

\node[raiz] (r) at (0,2.1) {Regla de puntuación propia\\(pérdida logarítmica, Brier)};
\node[caja] (cal) at (-1.75,0.75) {Calibración};
\node[caja] (ref) at ( 1.75,0.75) {Refinamiento};

\draw[flecha] (r.south) -- ++(0,-0.2) -| (cal.north);
\draw[flecha] (r.south) -- ++(0,-0.2) -| (ref.north);

\node[met] at (-1.75,-0.15) {¿El 70\,\% que promete corresponde al 70\,\% que ocurre?};
\node[met] at ( 1.75,-0.15) {¿Separa correctamente los casos que terminan distinto?};

\node[met, text=tecnavy] at (-1.75,-1.15) {mide: error de calibración, curva de fiabilidad};
\node[met, text=tecnavy] at ( 1.75,-1.15) {mide: área bajo la curva ROC};

\end{tikzpicture}
\caption{Descomposición de una regla de puntuación propia. Una métrica de
ordenamiento informa solo del componente derecho; una medida de calibración,
solo del izquierdo. Reportar únicamente una impide saber cuál de las dos
formas de fallar está ocurriendo.}
\label{fig:descomposicion}
\end{figure}

\begin{table}[!t]
\centering
\caption{Métricas de evaluación y pregunta que responde cada una.}
\label{tab:metricas}
\footnotesize
\begin{tabularx}{\columnwidth}{@{}P{1.6cm} P{1.3cm} Y@{}}
\toprule
\textbf{Métrica} & \textbf{Componente} & \textbf{Pregunta, y por qué no
basta por sí sola}\\
\midrule
Área bajo la curva ROC \newline \textit{(principal)} & Refinamiento &
¿Ordena correctamente dos posiciones? No depende de un umbral y es la
métrica del único antecedente comparable. No dice nada sobre si la
probabilidad es correcta\\
Pérdida logarítmica & Ambos &
¿Es exacta la probabilidad? Es además el criterio de ajuste
\eqref{eq:ajuste}, de modo que reportarla muestra si la optimización tuvo
éxito\\
Puntaje de Brier & Ambos &
¿Cuál es el error cuadrático de la probabilidad, y cómo se reparte entre
ambos componentes?\\
Curva de fiabilidad y error de calibración & Calibración &
Entre las posiciones a las que asigna un 70\,\%, ¿vence el 70\,\%?\\
Exactitud y medidas derivadas & --- &
Se reportan solo tras fijar y justificar un umbral, para comparabilidad con
trabajos que las emplean. Un modelo puede acertar la clase y estar
sistemáticamente sobreconfiado\\
\bottomrule
\end{tabularx}
\end{table}

El área bajo la curva es la métrica principal por tres razones: el desenlace
está cerca del equilibrio, de modo que la exactitud sería poco informativa;
la métrica no exige comprometerse con un umbral, y en esta etapa no hay
definido ningún costo de decisión; y es la única sobre la que existe una
cifra publicada análoga.

Las medidas de calibración se incluyen porque la evidencia indica que es la
propiedad que primero se degrada bajo cambio de población
\cite{guo_temporalshift_2021,xenopoulos_winprob_2022}, y porque ciertas
intervenciones sobre el balance de clases pueden dejar el ordenamiento
prácticamente intacto mientras distorsionan las probabilidades
\cite{vandenGoorbergh_imbalance_2022}.

Todas las métricas se reportan de forma global y por etapa de la partida.
Cada punto de esas curvas se acompaña del número de partidas y de
observaciones que lo sustenta, de modo que un desempeño reducido en la etapa
temprana pueda distinguirse de una simple escasez de datos en esa región.

Si la distribución de clases sobre la tabla final resultara marcadamente
desbalanceada, magnitud que depende de cuántas observaciones aporta cada
bando y que aún no se ha medido, se añadirá el área bajo la curva de
precisión y exhaustividad, declarando explícitamente cuál es la clase
minoritaria.

\noindent\textbf{Curvas de fiabilidad.} Los estimadores del error de
calibración basados en discretización dependen del número y la construcción
de los intervalos empleados \cite{silvafilho_calibration_2023}. Dimitriadis
\textit{et al.} proponen un procedimiento que evita esa dependencia y
produce diagramas reproducibles con una medida numérica asociada
\cite{dimitriadis_reliability_2021}. Se emplea si existe implementación
disponible en el entorno de trabajo; en caso contrario se usa el error de
calibración con discretización, declarando la limitación.

\subsubsection{Intervalos de confianza}

Todo intervalo se obtiene por remuestreo sobre partidas, nunca sobre filas,
por la misma razón que gobierna \eqref{eq:agrupacion}.

Los intervalos así construidos deben interpretarse como cotas optimistas de
la incertidumbre real. Brill \textit{et al.} reportan que el remuestreo que
ignora la agrupación alcanza una cobertura muy inferior a la nominal, y que
incluso el que la respeta permanece por debajo del nivel declarado
\cite{brill_winprob_2026}.

\subsection{Plan de modelos de referencia}
\label{subsec:planbaselines}

Cada peldaño de la escalera se justifica por lo que el anterior dejó sin
explicar.

\begin{table}[!t]
\centering
\caption{Escalera de modelos de referencia.}
\label{tab:baselines}
\footnotesize
\begin{tabularx}{\columnwidth}{@{}c P{2.1cm} Y@{}}
\toprule
\textbf{Nivel} & \textbf{Modelo} & \textbf{Pregunta que responde}\\
\midrule
B0 & Clasificador de probabilidad a priori & ¿Cuál es el piso? Lo que quede
por debajo es ruido\\
B1 & Heurística de dominio sobre una diferencia de recursos, fijada tras el
análisis exploratorio & ¿Hace falta aprendizaje automático?\\
B2 & Regresión logística regularizada & ¿Es la relación esencialmente
lineal? Aporta coeficientes inspeccionables\\
B3 & Árbol de decisión y bosques aleatorios & ¿Existen interacciones que el
modelo lineal no captura?\\
B4 & Potenciación del gradiente por histogramas & ¿Cuánto añade la
complejidad, y a qué costo de interpretabilidad?\\
\bottomrule
\end{tabularx}
\end{table}

La inclusión de B2 a B4 se apoya en la evidencia de que los conjuntos de
árboles constituyen el estado del arte sobre datos tabulares
\cite{borisov_tabular_2024,grinsztajn_trees_2022}.

\noindent\textbf{Sobre B1.} Hodge \textit{et al.} reportan que un modelo
construido sobre una única característica alcanza un desempeño cercano al
del conjunto completo \cite{hodge_winprediction_2021}. Un resultado análogo
aquí constituiría un hallazgo y no un fracaso, y declararlo de antemano evita
que, ante ese escenario, se busque justificar la complejidad del modelo con
argumentos posteriores a los datos.

\noindent\textbf{Sobre B2.} No se incluye por formalidad. Es el modelo que la
evidencia revisada trata como referencia honesta: fue el algoritmo más común
en los estudios de cambio temporal revisados \cite{guo_temporalshift_2021},
quedó a 0.016 de acierto del mejor modelo en el corpus análogo de
repeticiones \cite{bialecki_sc2egset_2023}, y fue de los menos inflados por
un protocolo defectuoso \cite{karbalaie_participantaware_2026}.

\noindent\textbf{Sobre B4.} Se emplea la implementación por histogramas
disponible en la biblioteca de aprendizaje automático que sostiene el resto
de la escalera. La elección atiende a dos criterios. Maneja de forma nativa
los valores ausentes, que en este corpus abundan por razones estructurales
descritas en la Sección~\ref{subsec:imputacion}. Y comparte interfaz con B2 y
B3, de modo que los cinco niveles se ejecutan bajo el mismo procedimiento sin
código específico por modelo, lo que reduce la superficie donde puede
introducirse una inconsistencia entre ejecuciones.

\subsubsection{Ablaciones por familia de características}

Las ablaciones eliminan un bloque completo del diccionario, progreso,
recursos, tablero, energía, estados alterados, diferenciales o contexto de
mazo, y reejecutan el mismo protocolo, de modo que el aporte marginal de
cada bloque se mida en lugar de afirmarse.

Una ganancia grande del bloque de contexto de mazo se interpreta como
evidencia de memorización del conjunto de estrategias dominante y no de
comprensión del estado: una representación que identifica el contenido en
lugar de describirlo puede ser excelente dentro de la distribución e inútil
fuera de ella \cite{bertram_cardrepresentations_2024}.

\subsubsection{Recalibración}
\label{subsec:recalibracion}

La corrección de probabilidades no se aplica de forma uniforme.
Niculescu-Mizil y Caruana reportan que los métodos de máximo margen presentan
distorsión sistemática, mientras que la regresión logística resulta bien
calibrada sin corrección y puede deteriorarse si se la recalibra
\cite{niculescumizil_probabilities_2005}. El procedimiento mide por tanto la
calibración de cada modelo antes de decidir si corregirla, y reporta ambas
mediciones, porque los mismos autores muestran que el ordenamiento de los
modelos según la calidad de sus probabilidades difiere antes y después de la
corrección.

Se adopta la regresión isotónica, siguiendo el criterio de tamaño de conjunto
de calibración que esos autores establecen: por encima del orden del millar
de casos, la isotónica iguala o supera a la alternativa paramétrica, y el
subconjunto reservado en la partición $60/15/10/15$ excede holgadamente ese
umbral. Ojeda \textit{et al.} recomiendan métodos basados en regresión
cuando el objetivo es transferir a poblaciones externas
\cite{ojeda_calibrating_2023}; dado que el escenario B constituye
precisamente ese caso, se reporta la calibración bajo ambos escenarios y se
contrasta si su recomendación aplica aquí.

\subsection{Plan de explicabilidad}
\label{subsec:explicabilidad}

La contribución no puede ser una puntuación aislada: la pregunta sustantiva
es cuáles características del estado transportan información sobre la
victoria. Se emplean cuatro instrumentos, cada uno ligado a una pregunta
distinta.

Los \textbf{coeficientes de la regresión logística} responden si la
asociación de una característica con la victoria tiene la dirección que el
conocimiento del juego predice.

La \textbf{importancia por permutación agrupada}, permutando dentro de cada
partida y nunca entre partidas, responde cuánto aporta una característica
una vez respetada la dependencia. La agrupación forma parte de la medición y
no es un detalle de ella, porque las estimaciones de importancia dependen del
diseño de validación y resultan inestables bajo particiones por fila
\cite{karbalaie_participantaware_2026}.

Las \textbf{ablaciones por familia} responden cuánto vale cada bloque del
diccionario.

Los \textbf{valores de atribución por instancia} se emplean solo cuando el
modelo seleccionado es un conjunto de árboles y la pregunta es por qué una
posición concreta recibió una puntuación concreta, que es lo que los tres
instrumentos anteriores no pueden responder.

\subsection{Contrato de reproducibilidad}
\label{subsec:reproducibilidad}

Las semillas y configuraciones se mantienen en ficheros versionados y no en
argumentos de línea de comandos; las dependencias quedan fijadas con
versiones explícitas; cada ejecución se registra junto con el identificador
de la revisión del código que la produjo; y se almacenan tablas comparables
de parámetros y métricas.

Los datos crudos no se versionan: se distribuyen las instrucciones y el
código para obtenerlos y reconstruir la tabla. Esto sigue la práctica del
descriptor de corpus análogo, que publica sus herramientas de extracción y
sus registros de procesamiento junto a los datos
\cite{bialecki_sc2egset_2023}, y atiende una debilidad documentada de la
literatura sobre cambio temporal, donde el código estaba disponible en solo
dos de quince estudios revisados \cite{guo_temporalshift_2021}.

Cada experimento se repite sobre múltiples semillas y se reporta la
dispersión resultante, práctica que los tres trabajos metodológicos de
referencia adoptan por razones equivalentes
\cite{grinsztajn_trees_2022,figueiredo_leakage_2024,brill_winprob_2026}.

\subsection{Riesgos y mitigación}
\label{subsec:riesgos}

\begin{table}[!t]
\centering
\caption{Riesgos identificados y su mitigación.}
\label{tab:riesgos}
\footnotesize
\begin{tabularx}{\columnwidth}{@{}P{1.8cm} c Y@{}}
\toprule
\textbf{Riesgo} & \textbf{Impacto} & \textbf{Mitigación}\\
\midrule
Observaciones de una misma partida repartidas entre subconjuntos & Alto &
Agrupación por partida en toda partición y todo pliegue; prueba automática
de intersección vacía\\[2pt]
Un campo excluido alcanza el vector de características & Alto &
Lista explícita; la construcción de la tabla falla si una columna prohibida
está presente\\[2pt]
Preprocesamiento ajustado sobre el conjunto completo & Alto &
Todas las transformaciones se ajustan dentro del pliegue de entrenamiento
\cite{sasse_leakage_2025}\\[2pt]
Cartas ausentes del entrenamiento tras particionar & Medio &
Auditoría de cobertura ejecutada después de particionar y antes de
entrenar\\[2pt]
Fuga por decisiones de diseño & Medio &
Reglas de exclusión y ventanas de turno declaradas antes de tocar el
conjunto de prueba \cite{kaufman_leakage_2011}\\[2pt]
Cambio de población entre días tempranos y tardíos & Medio &
Escenario B reportado por separado; cualquier caída se declara como
resultado y no como fallo \cite{guo_temporalshift_2021}\\[2pt]
Memorización de mazos en lugar de lectura del estado & Medio &
Escenario C y ablación del bloque de contexto de mazo
\cite{bertram_cardrepresentations_2024}\\[2pt]
Versión del motor confundida con el tiempo & Medio &
Registrar la versión como covariable de auditoría y declarar la
confusión\\[2pt]
Lectura optimista de una probabilidad no calibrada & Medio &
Reportar calidad probabilística junto a la capacidad de ordenar\\
\bottomrule
\end{tabularx}
\end{table}

\noindent\textbf{Expectativa declarada sobre el efecto del protocolo.} Se
anticipa que la partición agrupada produzca métricas inferiores y mayor
dispersión entre ejecuciones que una partición por fila. Ese comportamiento
está documentado \cite{figueiredo_leakage_2024} y su magnitud cuantificada
sobre datos de estructura comparable \cite{karbalaie_participantaware_2026},
y corresponde al funcionamiento correcto del protocolo y no a un defecto de
su ajuste. Declararlo antes de disponer de resultados protege la decisión de
mantener el protocolo ante métricas inferiores a las esperadas.

\subsection{Integración prevista con el agente}
\label{subsec:agente}

En la etapa final del proyecto, el estimador entrenado es consumido por un
agente inteligente como evaluador de posiciones. La interfaz es
deliberadamente estrecha: el agente entrega un estado, el mismo extractor de
características empleado en entrenamiento produce $x$, y el modelo devuelve
$f(x) \in [0,1]$.

El agente no entrena el modelo, no lo actualiza mediante interacción y no lo
sustituye, de modo que el proyecto no se convierte en aprendizaje por
refuerzo en ninguna etapa. El consumo se verifica mostrando que el
ordenamiento de posiciones candidatas por parte del agente cambia cuando el
estimador se reemplaza por el modelo trivial B0.

Esa integración es la razón última por la que la calibración importa: un
evaluador que ordena bien las posiciones pero está sistemáticamente
sobreconfiado engañará a cualquier regla de decisión que compare sus salidas
contra un umbral fijo.